% From the selected finite-volume mission to its limiting local price.
\section{From finite-volume mileage to the selected chart price}
\label{sec:ex-lattice-trace-mission-price}

This section proves the passage from the selected finite-volume mission to
its local price in the limiting table. The finite-volume separation estimate
of Proposition~\ref{prop:ex-lattice-finite-volume-pairing} supplies the
numerical bound used by the mileage argument. The mission's funded probe,
its local approximations, and the selected chart convergence are stated
explicitly below; this section does not claim to produce those analytic
receipts.

Write $i$ for the scale parameter, $j_i$ for its selected frame level,
and $k_i=\operatorname{summit}(j_i)$ for the associated volume scale. Let $a_i>0$ be the field scale,
$b_i>0$ the zoom factor, and $r_i=a_i b_i$ the physical insertion scale.
Let $\mathcal L$ be the ultrafilter on finite-volume indices $p=(N,s)$ used
by the construction, and define
\[
 \tau_i(q)=\exp\!\left(\frac{q}{10^4}
      (\log\widehat\beta_{j_i})^{48}\right).
\]
A mission consists of arrays $t_i(p)$ and $x_i$ such that, for all
sufficiently large $i$, $t_i(p)\to x_i$ along $\mathcal L$. Moreover,
there is one fixed integer $m$ such that for all sufficiently large $i$ and
for $\mathcal L$-almost every $p=(N,s)$, a probe of complexity at most
$m$, based at level $j_i$, has marked grade-five field $a_i^5t_i(p)$.
For such a probe, its measured marker $P$ has coefficient size at most
$a_i^5\widehat\beta_{k_i}^{m}$ and is odd under the central signed
inversion. Its toll receipt supplies a bounded local observable $U$ whose
pullback to the finite-volume chamber is interior and whose root error
relative to $P$ is at most $\tau_i(-200)$. For every pair of distinct
coarse sites $x,z$, it also supplies the bound
\[
 |\mathbb E_{i,p}[\widetilde U(x)\widetilde g(z)]|\leq\tau_i(-200),
\]
where tildes denote pullback to the finite-volume chamber. Here $g$ is any
coarse test whose pointwise absolute value is at most $\tau_i(150000)$,
which is interior to the fixed grade-four locality window, and whose
pullback is chamber-interior, provided that the window fits in the
available room at frame level $j_i$ inside the volume at scale $k_i$. For
the selected grade-four insertion, a localization receipt supplies one
coarse observable $g$ with all of the following properties: its pointwise
absolute value is at most $\tau_i(150000)$; it is interior to the fixed
grade-four locality window, and that window fits in the available room at
frame level $j_i$ inside the volume at scale $k_i$; its pullback
$V=\widetilde g$ is chamber-interior; the root error of $V$ relative to
the measured insertion is at most $\tau_i(-300)$; and the measured
insertion has root cost at most $\tau_i(5)$. The first three properties
are exactly the admissibility conditions on the test in the probe's
separated-site bound, so that bound applies to the pair
$(\widetilde U,V)$ with this particular $g$; a root approximation alone
would not permit this. The measured marker $P$ has root cost at most
$\tau_i(5)$ by a separate coefficient-to-root estimate, which uses its
coefficient bound and its oddness. These are the quantitative inputs
consumed by the finite-volume mileage estimate; their construction remains
an upstream obligation.

Fix a reference grade-four field $w$. Let $H_{i,p}(w)$ be its finite-height
insertion approximant, which is a raw grade-four label. It is determined
by the fixed reference lift of $w$ to a raw grade-four label, by the
grade-four transfer that the $i$th selected chart uses at the
finite-volume index $p$, by the zoom factor $b_i$ and by the frame level
$j_i$; its formula is not needed here. The mileage estimate is formulated
for insertions of exactly this form, under the conditions that this
transfer is a tested transfer at level $j_i$ and that
$1\le b_i\le\operatorname{tiltcap}(j_i)$; both conditions are part of the
selected chart receipt for large $i$. The selected
hierarchy receipt asserts $H_{i,p}(w)\to y_i(w)$ along $\mathcal L$, where
$y_i(w)$ is the limiting raw field and $\pi_4 y_i(w)=U_i^{(4)}w$. It also
ensures, on an
$\mathcal L$-large set, that the finite transfer is admissible, the marker
and insertion lie in the same reflection chamber, and the two coarse sites
are distinct. Let $\xi_i$ be the cover vector with first component $\lceil b_i\rceil$
and every other component $1$. Set $z_i=b_i^{-1}\xi_i$ componentwise.
The scale relation $r_i=a_i b_i$ gives $r_i z_i=a_i\xi_i$. Thus the
physical measurement point has the two cover coordinates
\[
 \ell_i=(\xi_i,0).
\]
The separation receipt for these coordinates supplies the distinct-site
hypothesis.

Set $\Theta_i=\tau_i(-100)$.
\begin{proposition}[Mission-to-chart price transfer]
\label{prop:ex-trace-mission-price-transfer}
Fix $w$. Assume the mission properties above. Assume that for all
sufficiently large $i$ and $\mathcal L$-almost every finite-volume index
$p$, mileage applies to the probe and insertion and gives
\begin{equation}
 \left|\mathbb E_{i,p}\!\left[
   A_{i,p}(a_i^5t_i(p);x_i^*)\,B_{i,p}(H_{i,p}(w);y_i^*)
                 \right]\right|\leq\Theta_i.
 \label{eq:ex-trace-mission-finite-price}
\end{equation}
Here $\mathbb E_{i,p}$ is the normalized finite-volume expectation,
$A_{i,p}$ and $B_{i,p}$ are the measured grade-five and grade-four
observables at the two distinct sites $x_i^*,y_i^*$ determined by $\ell_i$.
Assume the selected-table limit theorem:
whenever $u_p\to u$ and $v_p\to v$ along $\mathcal L$,
\begin{equation}
 \mathbb E_{i,p}[A_{i,p}(u_p;x_i^*)B_{i,p}(v_p;y_i^*)]
   \longrightarrow \mathcal M_i(\pi_5u,\pi_4v),
 \label{eq:ex-trace-mission-table-limit}
\end{equation}
where $\pi_d$ is the limiting projection in grade $d$ and $\mathcal M_i$
is the selected table's two-entry pairing at the scale-$r_i$ point.
Then, for all sufficiently large $i$ (with threshold allowed to depend on
$w$),
\[
 \left|\mathcal M_i\!\left(a_i^5\pi_5x_i,\ U_i^{(4)}w\right)\right|
      \leq\Theta_i.
\]
\end{proposition}
\begin{proof}
Fix a sufficiently large $i$ for which the mission, mileage, and chart
receipts hold. The set on which the probe and finite-volume insertion
conditions hold belongs to $\mathcal L$. On this set,
\eqref{eq:ex-trace-mission-finite-price} places every finite expectation in
the closed disk of radius $\Theta_i$. Mission convergence gives
$a_i^5t_i(p)\to a_i^5x_i$, and the chart receipt gives
$H_{i,p}(w)\to y_i(w)$ with $\pi_4y_i(w)=U_i^{(4)}w$. Applying
\eqref{eq:ex-trace-mission-table-limit} shows that the finite expectations
converge to $\mathcal M_i(a_i^5\pi_5x_i,U_i^{(4)}w)$; linearity of $\pi_5$
was used to take the fixed scalar $a_i^5$ outside. Since the disk is
closed, the limit satisfies the same bound. This proves the proposition.
\end{proof}

\subsection{The finite-volume mileage receipt}

The error scale is multiplicative:
$\tau_i(q)\tau_i(q')=\tau_i(q+q')$, and it is increasing in $q$. In the
finite-volume estimate, both measured markers have root cost at most
$\tau_i(5)$. The probe approximation contributes root error
$\tau_i(-200)$ and separated-site price $\tau_i(-200)$; the local
insertion approximation contributes root error $\tau_i(-300)$. The
reflection-positive product bound controls the two cross terms by products
of root costs. Bilinearity gives the exact decomposition
\[
 \langle P,Q\rangle=\langle U,V\rangle
   +\langle P,Q-V\rangle+\langle P-U,V\rangle.
\]
The three terms are bounded respectively by
\[
 \tau_i(-200),\qquad \tau_i(5)\tau_i(-300)=\tau_i(-295),\qquad
 \tau_i(-200)(\tau_i(5)+\tau_i(-300))
      =\tau_i(-195)+\tau_i(-500).
\]
Thus the sum is at most $4\tau_i(-195)$. The selected growth receipt gives
$4\leq\tau_i(1)$ eventually, hence
$4\tau_i(-195)\leq\tau_i(-194)\leq\tau_i(-100)$.
Proposition~\ref{prop:ex-lattice-finite-volume-pairing} proves this
three-term estimate from the explicit root and reflection-positive bounds.
The remaining hard content is construction of the probe, local
approximation, and root-cost receipts at the selected scales.

In the source route, the finite-cover construction produces the mission's
fixed-complexity probe, while convergence of the selected core field
produces its limiting value. The limit transfer above is complete once
those receipts are supplied. Producing the funded cover, the probe witness,
and the quantitative localization receipts remains part of the larger
construction.
